% Fragmento público generado desde propietarios íntegros.
% Residencia: 52.4.1.
% Fuente: ../incoming/radion_monografia_20260827/manuscrito/sections/05_geometria_accion.tex:1-217
\noindent Spectral geometry determines the form, and the areal law determines the action transported by the phase.\par\medskip

\subsubsection{Equivalence between the spectral and geometric charts}

\paragraph{Spectral positive operator and conjugate eigenvalues}

The angular momentum produces the two-layer operator.
\begin{equation}
B_1=AI+C^*R,\qquad R^2=I,\qquad
P_\pm=\frac{I\pm R}{2}.
\label{eq:B1}
\end{equation}
Its eigenvalues are \(A\pm C^*\). Since \(0<C^*<A\), both are positive.
In the chart \(\kappa=\pi/180\), the spectral ellipse has semiaxes
\begin{equation}
a_1^2=\kappa(A+C^*),\qquad
b_1^2=\kappa(A-C^*).
\label{eq:semiejes-espectrales}
\end{equation}
The decomposition is reversible:
\begin{equation}
A=\frac{a_1^2+b_1^2}{2\kappa},\qquad
C^*=\frac{a_1^2-b_1^2}{2\kappa}.
\label{eq:reconstruccion-AC}
\end{equation}
The ellipse is not drawn to illustrate a number. It exactly reconstructs the pair that generated it.

Define
\begin{equation}
\rho_1=\sqrt{A^2-C^{*2}},\qquad
\eta=\operatorname{artanh}\frac{C^*}{A}
=\frac12\log\frac{A+C^*}{A-C^*}.
\label{eq:eta}
\end{equation}
Then
\[
A\pm C^*=\rho_1e^{\pm\eta},
\]
and
\begin{equation}
a_1=\sqrt{\kappa\rho_1}\,e^{\eta/2},
\qquad
b_1=\sqrt{\kappa\rho_1}\,e^{-\eta/2}.
\label{eq:semiejes-eta}
\end{equation}
The product preserves scale; the quotient preserves form:
\[
a_1b_1=\kappa\rho_1,
\qquad
\frac{b_1}{a_1}=e^{-\eta}
=\sqrt{\frac{A-C^*}{A+C^*}}.
\]

\paragraph{Focal Invariants of the Positive Operator}

The semifocal distance is
\begin{equation}
c_1^2=a_1^2-b_1^2=2\kappa C^*,
\qquad
F_\pm=(\pm c_1,0).
\label{eq:focos}
\end{equation}
The total distance is equal to
\begin{equation}
d_1=2c_1=2\sqrt{2\kappa C^*}
=0.544545509325626623406299779866023\ldots.
\label{eq:d1}
\end{equation}
The isolated number is not the thesis. The identity \(c_1^2=2\kappa C^*\) states that
\(C^*\) is literally the quadratic focal splitting in the
distinguished chart. The foci are the conjugate endpoints of the spectral aperture;
they are not, by themselves, the electric and magnetic poles. That reading requires
the constitutive transducer.

The eccentricity is
\begin{equation}
e_f=\frac{c_1}{a_1},\qquad
e_f^2=1-e^{-2\eta}=\frac{2C^*}{A+C^*}.
\label{eq:excentricidad}
\end{equation}
Numerically,
\[
\begin{aligned}
\eta&=0.299689683921630799684926562863927\ldots,\\
e_f&=0.671451895550191986916277055864332\ldots.
\end{aligned}
\]

\paragraph{Reference APP Block and Normalized Transport}

The central block
\[
B_0=7I+2R=9P_++5P_-
\]
possesses
\[
\rho_0=\sqrt{45},\qquad
\eta_0=\frac12\log\frac95,\qquad
e_0=\frac23.
\]
Its focal length is
\[
d_0=4\sqrt\kappa
=0.528443639680801468446003835349358\ldots.
\]
The exact comparison is
\begin{equation}
\boxed{\left(\frac{d_1}{d_0}\right)^2=\frac{C^*}{2}},
\qquad
d_1=d_0\sqrt{\frac{C^*}{2}}.
\label{eq:invariante-focal-relativo}
\end{equation}
This is the defensible genealogy of \(0.544545\ldots\): a focal APP scale
of reference multiplied by a relative invariant. Its visual proximity to
\(0.54\) is not an identity with the APP defect \(54\).

Transport from \(B_0\) separates into scale and anisotropy:
\begin{equation}
T_{\mathrm{rel}}=\frac{\rho_1}{\sqrt{45}}
\exp\bigl((\eta-\eta_0)R\bigr).
\label{eq:transporte-rel}
\end{equation}
On the leaves,
\[
t_+=\frac{A+C^*}{9}=1.0467878786947163\ldots,
\qquad
t_-=\frac{A-C^*}{5}=1.0347228460630311\ldots.
\]
The deformation is not a single decimal. Its isotropic component is \(\sqrt{t_+t_-}\); its hyperbolic component is \(\tfrac12\log(t_+/t_-)\).

\subsubsection{Positive ellipse associated with the angular pair}

\paragraph{Construction via Conjugate Semi-axes}

The preceding spectral scale has no units of action. Dimensional
Mechanics publishes the same form on the internal action line:
\begin{equation}
a_S=\sqrt{2\hret e^\eta},
\qquad
b_S=\sqrt{2\hret e^{-\eta}}.
\label{eq:semiejes-accion}
\end{equation}
Immediately,
\begin{equation}
a_Sb_S=2\hret,\qquad
\frac{a_S}{b_S}=e^\eta.
\label{eq:producto-forma}
\end{equation}
The product fixes the action; the quotient fixes the anisotropy. The ellipse is
parametrized by
\begin{equation}
Q(\theta)=a_S\cos\theta,\qquad
P(\theta)=b_S\sin\theta.
\label{eq:elipse-param}
\end{equation}

The spectral and action ellipses are homothetic:
\begin{equation}
\lambda_{\mathrm{act}}^2=\frac{2\hret}{\kappa\rho_1},
\qquad
(a_S,b_S,c_S)=\lambda_{\mathrm{act}}(a_1,b_1,c_1).
\label{eq:homotecia}
\end{equation}
They share the form \(e^{-\eta}\), not the physical dimension.

\paragraph{Areal law of action per radion}

The forma simplectica canonica is \(\omega=\diff Q\wedge\diff P\). The area
oriented barrida from \(0\) until \(\theta\) is
\begin{align}
\mathcal A(\theta)
&=\frac12\int_0^\theta
\bigl(Q(t)P'(t)-P(t)Q'(t)\bigr)\,\diff t\\
&=\frac12a_Sb_S\int_0^\theta
(\cos^2t+\sin^2t)\,\diff t.
\end{align}
Using \eqref{eq:producto-forma}, we obtain the central result.

\begin{theorem}[Action per radion]
For the ellipse \eqref{eq:elipse-param},
\begin{equation}
\boxed{\mathcal A(\theta)=\hret\theta.}
\label{eq:area-theorem}
\end{equation}
In particular,
\begin{equation}
\boxed{\mathcal A(1)=\hret},
\qquad
\boxed{\mathcal A(2\pi)=2\pi\hret=\hfull}.
\label{eq:area-radion-vuelta}
\end{equation}
\end{theorem}

This equality allows the definition of the radius without first appealing to the arc length.

\begin{definition}[Oriented radion]
\begin{equation}
\mathfrak r_\sigma=(\theta=1,\sigma),
\qquad \sigma\in\{+1,-1\},
\qquad
\mathcal A(\mathfrak r_\sigma)=\sigma\hret.
\label{eq:radion-oriented}
\end{equation}
\end{definition}

The radian is the projection that forgets \(\sigma\), the sheet, the
memory, and the area. The radion preserves that information. Hence the name
\emph{rad--ion}: orientation precedes
physical charge, and charge appears later as
\[
Q(x)=n_Q(x)e_{\mathrm{int}},\qquad
e_{\mathrm{int}}=\sqrt{\frac{2\alphah\hfull}{Z_{\mathrm{int}}}}.
\]

